\documentclass[12pt,a4paper]{article}
\usepackage{iftex}
\ifPDFTeX
    \usepackage[utf8]{inputenc}
    \usepackage[T1]{fontenc}
    \usepackage{mathpazo}
\else
    \usepackage{fontspec}
    \setmainfont{Times New Roman}
\fi
\usepackage[brazilian]{babel}
\usepackage{amsmath, amssymb, amsthm}
\usepackage[linesnumbered,ruled,vlined,portuguese]{algorithm2e}
\usepackage{listings}
\usepackage{booktabs}
\usepackage{geometry}
\usepackage{float}
\usepackage{xcolor}
\usepackage{setspace}
\usepackage[colorlinks=true, linkcolor=blue, urlcolor=blue, citecolor=blue]{hyperref}

\geometry{margin=2.5cm}
\onehalfspacing

\definecolor{codebg}{rgb}{0.96,0.96,0.96}
\definecolor{codekey}{rgb}{0.0,0.4,0.8}
\definecolor{codestr}{rgb}{0.6,0.1,0.1}
\definecolor{codecom}{rgb}{0.4,0.6,0.4}

\lstset{
    backgroundcolor=\color{codebg},
    commentstyle=\itshape\color{codecom},
    keywordstyle=\bfseries\color{codekey},
    stringstyle=\color{codestr},
    basicstyle=\ttfamily\footnotesize,
    breaklines=true,
    numbers=left,
    numberstyle=\tiny\color{gray},
    frame=single,
    rulecolor=\color{lightgray},
    captionpos=b,
    tabsize=4
}

\title{
    \vspace{-1cm}
    \begin{center}
        \Large \textbf{Universidade Federal do Ceará -- UFC} \\
        \large Departamento de Computação \\
        \vspace{1cm}
        \Huge \textbf{Métodos Numéricos II} \\
        \vspace{0.5cm}
        \Large \textit{Relatório da Unidade 3: Autovalores e SVD}
    \end{center}
}
\author{
    Nicolas Wallace Amorim Perote -- Matrícula: 566512 \\
    Samyra Vitória Lima de Almeida -- Matrícula: 521240
}
\date{\today}

\begin{document}

\maketitle
\tableofcontents
\newpage

\section{Introdução}
O cálculo algébrico-numérico de autovalores ($\lambda$) e autovetores ($v$) de matrizes quadradas é o pilar estrutural para inumeráveis áreas, da física quântica à engenharia de telecomunicações (ex: Compressão via SVD e Análise de Componentes Principais). Resolver analiticamente a equação $\det(A - \lambda I) = 0$ requer encontrar raízes de polinômios de grau $n$, o que se torna matematicamente inviável (Pelo Teorema de Abel-Ruffini) ou ineficiente para $n \geq 5$.

O escopo da \textbf{Unidade 3} concentra-se no delineamento algorítmico da abstração espectral de uma matriz. Foram propostas e implementadas suítes em \textit{Python} puramente vetorizadas utilizando a infraestrutura do \texttt{NumPy} para contornar gargalos presentes em métodos tradicionais estruturados por laços imperativos.

Os métodos transitam entre abordagens \textit{ad-hoc} focadas na extração de valores extremos (Métodos da Potência Regular e Inversa) até estratégias robustas de diagonalização baseadas em similaridade ortogonal (Householder, QR e Jacobi), culminando com a Decomposição em Valores Singulares (SVD).

\section{Métodos de Potência}
A família de métodos da potência explora propriedades projetivas do espectro de $A$.
Dado um vetor aleatório $v_0$, a multiplicação sucessiva $v_{k+1} = \frac{A v_k}{||A v_k||}$ alinha o vetor, inevitavelmente, com a direção correspondente ao autovalor de maior magnitude $|\lambda_{max}|$. Este procedimento constitui o \textbf{Método da Potência Regular}.

Para obter o menor autovalor $|\lambda_{min}|$, a álgebra prescreve aplicar a mesma técnica sobre $A^{-1}$. Computacionalmente, o \textbf{Método da Potência Inversa} abdica de inverter explicitamente a matriz. No nosso módulo unificado (\texttt{EigenMethods}), extraímos os fatores através de `scipy.linalg.lu\_factor` e resolvemos o sistema triangular iterativamente, convertendo a complexidade assintótica $\mathcal{O}(n^3)$ para $\mathcal{O}(n^2)$ por iteração.

Adicionando um parâmetro escalar ao sistema $(A - \mu I)^{-1}$, criamos o \textbf{Método da Potência com Deslocamento}, o qual permite sondar o espectro e extrair a raiz espectral rigorosamente mais próxima do palpite $\mu$.

\begin{lstlisting}[language=Python, caption={Extrato da Implementação da Potência Inversa com Fatoração LU}]
@classmethod
def power_method_inverse(cls, A, v0, tol=1e-6, max_iter=1000):
    v = cls._normalize(np.array(v0, dtype=float))
    lu, piv = scipy.linalg.lu_factor(A)
    lam_old = 0.0
    for _ in range(max_iter):
        y = scipy.linalg.lu_solve((lu, piv), v)
        v_new = cls._normalize(y)
        lam_new = v_new @ (A @ v_new)
        if np.abs(lam_new - lam_old) < tol:
            return lam_new, v_new
        lam_old = lam_new
        v = v_new
\end{lstlisting}

\section{Decomposições Ortogonais: Householder, QR e Jacobi}
Uma base inteiramente diferente para exploração espectral repousa nas transformações de similaridade. Se $Q$ é ortogonal ($Q^{-1} = Q^T$), a matriz $B = Q^T A Q$ detém o idêntico conjunto de autovalores que $A$.

\subsection{Reflexões de Householder}
Uma matriz reflexiva do tipo $H = I - 2nn^T$ (onde $n$ é unitário) espelha vetores por hiperplanos. Por intermédio da biblioteca \texttt{EigenMethods.householder\_tridiagonalize}, aplicamos sequenciais reflexões para sistematicamente aniquilar elementos distantes da diagonal matriz $A$, originando um núcleo \textit{tridiagonal simétrico}. 
Isto acelera formidavelmente o custo de algoritmos subsequentes, como o QR.

\subsection{O Algoritmo QR Iterativo}
Pela Fatoração QR, a matriz é segregada em $A = QR$. Na versão de autovalores, recriamos $A_{k+1} = R Q$.
Esta reversão sutil efetua a projeção $Q^T A Q$, eliminando vagarosamente os resíduos fora-diagonais. Desenvolvemos o algoritmo munido de limiares numéricos rigorosos baseados na restrição Euclidiana. Ao aplicar à matriz $A_3$ (fornecida para testes), verificou-se a matriz convergir formidavelmente, expondo seu espectro completo: $\{49.38, 31.31, 23.64, 11.64, 4.01\}$.

\subsection{Diagonalização de Jacobi}
Enquanto o QR ataca $A$ recursivamente, o \textbf{Método de Jacobi} utiliza matrizes de Rotação (Givens) puras, focando analiticamente no aniquilamento pontual do maior elemento remanescente $a_{ij}$. 
\begin{equation}
\theta = \frac{1}{2}\arctan\left(\frac{2a_{ij}}{a_{ii}-a_{jj}}\right)
\end{equation}
Por atuar diretamente sobre os maiores resíduos (monitorado no algoritmo usando \texttt{np.argmax}), garantiu formidável robustez e convergiu iterativamente para os mesmos autovalores descobertos via algoritmo QR.

\section{Decomposição em Valores Singulares (SVD)}
Qualquer matriz arbitrária de dimensões $(m \times n)$ admite a fatoração absoluta:
\begin{equation}
A = U \Sigma V^T
\end{equation}
Os vetores singulares esquerdos $U$ e direitos $V$ carregam as direções de variância, enquanto $\Sigma$ dita as escalas.
Na nossa classe especial \texttt{SVDMethods}, contornamos implementações caixa-preta.
Extraímos $\Sigma$ obtendo os autovalores $(\lambda_i)$ da gramiana normal $A A^T$ (se $m<n$) ou $A^T A$ (se $m \ge n$), executados perfeitamente pelo módulo de \textit{Jacobi} já estabelecido.
Os valores singulares provêm diretamente de $\sigma_i = \sqrt{\lambda_i}$. Ao remontar a matriz a partir do trio $(U, \Sigma, V^T)$ numa pequena malha real $2 \times 3$, reportamos resíduo espúrio da ordem do $\epsilon$ da máquina (na ordem de $10^{-16}$).

\section{Conclusão}
Os métodos investigados na Unidade 3 oferecem um vislumbre das arquiteturas \textit{Standard} (BLAS/LAPACK) adotadas globalmente pela computação moderna. A migração das rotinas antigas, lotadas de iteradores \texttt{for}, para um ecossistema Python O.O. limpo, pautado essencialmente pelo núcleo matricial transposto em C pelo \texttt{NumPy}, ratifica um enorme progresso estrutural.

De simples vetores convergindo por método da potência ao suntuoso balé das rotações unificadas de Jacobi construindo uma base SVD exata, as abstrações descritas no relatório denotam plenas capacidades teóricas e práticas em manipulação e inteligência matricial.

\end{document}
